(1) 取$H_y^{(d)}=A e^{ik_xx-\alpha_d z-i\omega t}$、$H_y^{(m)}=A e^{ik_xx+\alpha_m z-i\omega t}$，其中$\alpha_j=\sqrt{k_x^2-\epsilon_j k_0^2}$并选$\operatorname{Re}\alpha_j>0$。由$\nabla\times H=-i\omega\epsilon_0\epsilon_jE$得到$E_x^{(d)}\propto-\alpha_d H_y^{(d)}/\epsilon_d$、$E_x^{(m)}\propto+\alpha_mH_y^{(m)}/\epsilon_m$；切向$E_x,H_y$连续给$\alpha_d/\epsilon_d+\alpha_m/\epsilon_m=0$。
(2) 联立$\alpha_j^2=k_x^2-\epsilon_j k_0^2$及上述边界式，消元得$k_x^2=k_0^2\epsilon_d\epsilon_m/(\epsilon_d+\epsilon_m)$，选择正向传播支$k_x=k_0\sqrt{\epsilon_d\epsilon_m/(\epsilon_d+\epsilon_m)}$。纯实无损束缚模存在条件$\epsilon_m<-\epsilon_d$，此时$k_x>k_0\sqrt{\epsilon_d}$位于介质光锥之外；两侧穿透长度分别为$1/\alpha_d$和$1/\alpha_m$。
(3) Drude函数满足$\epsilon_m<-\epsilon_d$仅当$0<\omega<\omega_{\rm sp}=\omega_p/\sqrt{1+\epsilon_d}$。$\omega\to0$时$\epsilon_m\to-\infty$，故$k_x\sim k_0\sqrt{\epsilon_d}$，表面模接近介质光线；$\omega\uparrow\omega_{\rm sp}$时$\epsilon_m+\epsilon_d\to0^-$，传播常数$k_x\to\infty$，局域场尺度缩小，表面等离激元频率渐近$\omega_{\rm sp}$。
(4) 对非磁性两介质，TE模连续切向场会要求$\alpha_d+\alpha_m=0$，与两个正衰减常数矛盾，因此无对应束缚TE模。TM模式的时间平均能流须包含金属色散储能，不能把$\epsilon_m<0$直接当成负总能量；若考虑$\operatorname{Im}\epsilon_m>0$的吸收，则$k_x$成为复数，沿$x$的传播振幅衰减，无损$k_x\to\infty$的奇异极限被阻尼和非局域效应钝化。
